\subsection{Folded Regulated Cascode Op-Amp}\label{sec:amp_folded_regulated}

\subsubsection{Overview}
The folded cascode of Section~\ref{sec:amp_folded} with both cascodes
(\(M_3\) and \(M_5\)) gain boosted as in Section~\ref{sec:amp_regulated}
by auxiliary amplifiers \(A_{bn}\), \(A_{bp}\). Folding keeps the input
range, regulation recovers (and exceeds) the telescopic gain.

\subsubsection{Schematic (half circuit)}
\begin{circuitfig}
  \draw (0,0) node[nmos](MB){}; \node[right] at (0.3,0) {$M_B$};
  \draw (0,2.2) node[nmos](M3){}; \node[right] at (0.3,2.2) {$M_3$};
  \draw (0,4.5) node[pmos](M5){}; \node[right] at (0.3,4.5) {$M_5$};
  \draw (0,6.7) node[pmos](M7){}; \node[right] at (0.3,6.7) {$M_7$};
  \draw (MB.G) -- ++(-0.5,0) node[left]{$V_{BB}$};
  \draw (M7.G) -- ++(-0.5,0) node[left]{$V_{B7}$};
  \draw (MB.S) node[ground]{};
  \draw (MB.D) -- (M3.S);
  \draw (M3.D) -- (M5.D);
  \draw (M5.S) -- (M7.D);
  \draw (M7.S) node[vcc]{$V_{dd}$};
  \draw (0,1.1) node[circ]{} node[below right]{$V_f$};
  \draw (-2.2,2.0) node[op amp, scale=0.6, yscale=-1](BN){};
  \draw (0,1.1) -- (-3.6,1.1) |- (BN.-);
  \draw (BN.+) -- ++(-0.6,0) node[left]{$V_{refn}$};
  \draw (BN.out) -- ++(0.3,0) |- (M3.G);
  \node at (-2.2,2.8) {$A_{bn}$};
  \draw (0,5.6) node[circ]{} node[right]{$V_u$};
  \draw (-2.2,4.7) node[op amp, scale=0.6](BP){};
  \draw (0,5.6) -- (-3.6,5.6) |- (BP.-);
  \draw (BP.+) -- ++(-0.6,0) node[left]{$V_{refp}$};
  \draw (BP.out) -- ++(0.3,0) |- (M5.G);
  \node at (-2.2,3.9) {$A_{bp}$};
  \draw (0,3.35) node[circ]{} -- (1.4,3.35) node[circ]{} node[above]{$V_{o1}$};
  \draw (3,6.0) node[vcc]{$V_{dd}$} to[american current source, l={$\frac{I_T}{2} + \frac{g_{m1} V_{id}}{2}$}] (3,1.1) -- (0,1.1);
\end{circuitfig}

\subsubsection{Transfer Function}
\begin{enumerate}
  \item \textbf{Large Signal, Time Domain, DC Input Analysis (Q5)}: the
    currents are those of the folded cascode. The loops set \(V_F = V_{refn}\)
    and \(V_U = V_{refp}\); with \(V_{refn} = V_{ov,B}\) and
    \(V_{dd} - V_{refp} = |V_{ov7}|\) the swing is the folded cascode's,
    \(V_{O,\min} = V_{ov,B} + V_{ov3}\), \(V_{O,\max} = V_{dd} - |V_{ov7}| - |V_{ov5}|\).
  \item \textbf{Small Signal, Frequency Domain, AC Input Analysis (Q4)}: the
    references are fixed, so \(V_{g3} = A_{bn}(0 - V_f) = -A_{bn} V_f\) (the fold node
    sits at the inverting input of \(A_{bn}\)) and \(V_{g5} = -A_{bp} V_u\). Body
    effect is neglected (\(g_{mb} = 0\)), as in Section~\ref{sec:amp_regulated}.\\
    \textbf{Output Impedance Down} (\(I_{in} = 0\), test current \(I_t\) down through
    \(M_3\) and \(R_f = r_{o1} \parallel r_{oB}\), so \(V_f = I_t R_f\)):
    \begin{align*}
      I_t &= g_{m3}(V_{g3} - V_f) + g_{o3}(V_t - V_f) = g_{m3}(-A_{bn} V_f - V_f) + g_{o3} V_t - g_{o3} V_f \\
      I_t &= -\left(g_{m3}(1 + A_{bn}) + g_{o3}\right) I_t R_f + g_{o3} V_t \\
      I_t\left(1 + g_{m3}(1 + A_{bn}) R_f + g_{o3} R_f\right) &= g_{o3} V_t \\
      Z_{dn} &= \frac{V_t}{I_t} = r_{o3} + g_{m3}(1 + A_{bn}) r_{o3} R_f + R_f
    \end{align*}
    the folded cascode \(Z_{dn}\) with \(g_{m3} \to g_{m3}(1 + A_{bn})\).
    \textbf{Output Impedance Up}: the same steps for \(M_5\) over \(r_{o7}\):
    \(Z_{up} = r_{o5} + g_{m5}(1 + A_{bp}) r_{o5} r_{o7} + r_{o7}\).\\
    \textbf{Short Circuit Current} (\(V_{o1} = 0\)): \(I_{d3} = g_{m3}(-A_{bn} V_f - V_f) + g_{o3}(0 - V_f) = -\left(g_{m3}(1 + A_{bn}) + g_{o3}\right) V_f\);
    KCL at \(V_f\) as in Section~\ref{sec:amp_folded}:
    \begin{align*}
      |I_{sc}| &= \frac{g_{m3}(1 + A_{bn}) + g_{o3}}{g_{m3}(1 + A_{bn}) + g_{o3} + g_{o1} + g_{oB}} I_{in} \approx I_{in}
    \end{align*}
    \textbf{Gain}: with \(I_{in} = \frac{g_{m1} V_{id}}{2}\) and \(V_{o2} = -V_{o1}\) as before,
    \begin{align*}
      |A_d| &= \frac{2 |I_{sc}| (Z_{dn} \parallel Z_{up})}{V_{id}} \approx g_{m1}(Z_{dn} \parallel Z_{up})
             \approx g_{m1}\left[g_{m3} A_{bn} r_{o3} R_f \parallel g_{m5} A_{bp} r_{o5} r_{o7}\right]
    \end{align*}
    keeping the dominant term of each impedance (\(1 + A_b \approx A_b\)).
\end{enumerate}

\subsubsection{Analysis}
\begin{enumerate}
  \item \textbf{Gain Estimate}: take every \(g_m r_o\) equal, \(R_f \approx r_o\) and
    \(A_{bn} = A_{bp} = A_b\). Then
    \begin{align*}
      Z_{dn} \approx Z_{up} &\approx g_m r_o \cdot A_b \cdot r_o \hspace{1cm}
      Z_{dn} \parallel Z_{up} \approx \frac{g_m r_o A_b r_o}{2} \\
      |A_d| &\approx g_m \cdot \frac{g_m r_o A_b r_o}{2} = \frac{(g_m r_o)^2 A_b}{2}
    \end{align*}
    With single stage auxiliary amplifiers \(A_b \approx g_m r_o\):
    \(|A_d| \approx \frac{(g_m r_o)^3}{2}\); e.g. \(g_m r_o = 50\):
    \(\frac{50^3}{2} = \frac{125000}{2} = 62500\) (\(20\log_{10} 62500 = 96\,dB\)) in a single stage.
  \item \textbf{Gain-Bandwidth}: \(GB = \frac{|A_d|\, p_1}{2\pi}\) with
    \(p_1 = \frac{1}{(Z_{dn} \parallel Z_{up}) C_L}\) gives \(GB = \frac{g_{m1}}{2\pi C_L}\), unchanged.
    The auxiliary amplifiers must have their unity gain frequency between the
    main loop bandwidth and the main non-dominant pole to avoid slow doublets
    (Section~\ref{sec:amp_regulated}).
\end{enumerate}
